% ============================================
% metadata.tex
% Berisi definisi semua variabel untuk dokumen usulan penelitian S3
% ============================================

% Jenis dokumen
% \newcommand{\jenisdokumen}{USULAN PENELITIAN S3}
\newcommand{\jenisdokumen}{UJIAN KOMPREHENSIF}
% \newcommand{\jenisdokumen}{UJIAN KOMPREHENSIF}

% Judul penelitian
\newcommand{\judul}{\MakeUppercase{Adaptasi CDE-Mapper Pemetaan Terminologi Klinis untuk Teks Bahasa Indonesia: Pendekatan Berbasis LLM dan Hybrid Retrieval}}

% Penulis dan NIM
\newcommand{\penulis}{Anie Rose Irawati}
\newcommand{\nim}{24/551784/SPA/01084}

% Program studi, departemen, fakultas, universitas
\newcommand{\programstudi}{S3 Ilmu Komputer}
\newcommand{\departemen}{Departemen Ilmu Komputer dan Elektronika}
\newcommand{\fakultas}{Fakultas Matematika dan Ilmu Pengetahuan Alam}
\newcommand{\universitas}{Universitas Gadjah Mada}

% Tanggal sidang/pengesahan
\newcommand{\tanggal}{.......... 2026}

% Dewan penguji (gunakan placeholder jika belum diisi)
\newcommand{\ketuaPenguji}{...............}   % masih kosong
\newcommand{\promotor}{Afiahayati, S.Kom., M.Cs., Ph.D}
\newcommand{\kopromotor}{Dr. Lukman Heryawan, S.T., M.T.}
\newcommand{\pengujiSatu}{Prof Dr. XYZ, M.SC}   % placeholder
\newcommand{\pengujiDua}{Prof Dr. ABC, M.SC}    % placeholder
\newcommand{\pengujiTiga}{Prof Dr. XYZ, M.SC}   % placeholder (penguji tambahan)